
\subsubsection{2.1 Execution-Based Feedback for Code Generation}

Execution feedback uses unit test outcomes as training signals. CodeRL \citep{le2022coderl} applies actor-critic reinforcement learning for code, using test pass/fail as episode-level rewards. A critic network predicts functional correctness, enabling critical sampling for code regeneration.

RLTF \citep{liu2023rltf} extends execution feedback with multi-granularity signals. Beyond binary pass/fail, RLTF extracts fine-grained error locations from unit test output, providing line-level localization. Online RL with real-time data generation is reported to yield strong results on APPS and MBPP benchmarks.

PPOCoder \citep{shojaee2023ppocoder} combines PPO-based RL with non-differentiable execution feedback and structure alignment.

These methods share a characteristic: execution feedback provides no semantic guidance about why code fails.

\subsubsection{2.2 AI-Generated Feedback for Code}

AI feedback uses LLM-generated critique for code improvement. Self-Refine \citep{madaan2023selfrefine} demonstrates zero-shot iterative refinement: the same LLM generates code, critiques it, and refines based on its own feedback. No training is required---refinement occurs at inference time.

However, AI feedback has low fidelity. Madaan et al. analyze Self-Refine failures: 33\% stem from identifying the wrong error location, and 61\% from proposing incorrect fixes.

\subsubsection{2.3 Gap}

No prior work uses execution as a verification mechanism for AI feedback during training. Execution and AI feedback have been developed as alternatives rather than combined.
